% TEMPLATE-MILC-v3.tex - plantilla maestra UNICA de las guias MILC v3.
%
% La usan las tres familias de guias, cada una con su driver:
%   · Grados 6-11  -> scripts/build-guias-g11.py
%   · Bebras       -> scripts/build-guias-bebras.py
%   · Semillero    -> scripts/build-guias-semillero.py
%
% Antes habia tres copias casi identicas (~400 lineas iguales, 6 distintas):
% cada mejora habia que aplicarla tres veces, y en la practica no se hacia
% (el motor de recursos vivia solo en dos de ellas). Lo que varia entre
% familias esta parametrizado en 6 placeholders que cada driver llena
% (se nombran aqui SIN su sintaxis real: el builder reemplaza por texto plano,
% tambien dentro de los comentarios, y un valor multilinea rompe el preambulo):
%
%   HEADER_IZQ        encabezado izquierdo de pagina
%   HEADER_DER        encabezado derecho de pagina
%   PORTADA_ETIQUETA  rotulo sobre el numero grande (GRADO/MOMENTO/MODULO)
%   PORTADA_SERIE     linea de serie de la portada
%   PORTADA_ID        identificador de la guia en la portada
%   PORTADA_CIRCULO   el circulito de grado (°), o vacio si no aplica
%
% El resto de claves las documenta content/guias/_SCHEMA.md. El script valida
% que no queden placeholders sin reemplazar antes de escribir el archivo final.

\documentclass[11pt,letterpaper]{article}
\usepackage[margin=2.0cm]{geometry}
\usepackage[table]{xcolor}
\usepackage{fontspec}
\usepackage[spanish]{babel}
\usepackage{tikz}
% Librerías TikZ para los diagramas de las guías (bloque `recursos:`):
%  · babel  → babel en español vuelve activos caracteres como «>», y TikZ los usa
%             en sus opciones (>=stealth, ->). Sin esta librería, cualquier
%             diagrama con flechas revienta con «\language@active@arg> has an
%             extra }». Debe ir ANTES de que se dibuje nada.
%  · positioning        → sintaxis «below left=2pt and 3pt».
%  · decorations.pathreplacing → llaves ({brace}) para señalar rangos.
\usetikzlibrary{babel,positioning,decorations.pathreplacing}
\usepackage{graphicx}
% Circuitos: el trabajo de electrónica se hace hoy con micro:bit y sus sensores
% integrados (MakeCode, sin cableado), así que no se necesitan esquemáticos.
% Si algún día entran montajes con protoboard, basta descomentar esta línea:
% los diagramas .tex/.tikz declarados en `recursos:` ya se incrustan solos.
% \usepackage{circuitikz}
\usepackage[most]{tcolorbox}
\usepackage{tabularx}
\usepackage{array}
\usepackage{fancyhdr}
\usepackage{titlesec}
\usepackage[document]{ragged2e}  % bandera a la derecha en todo el cuerpo
\usepackage{enumitem}
\usepackage{microtype}

% Helvetica Neue no trae la flecha U+2192, asi que cada «→» escrito literalmente
% en el YAML se compilaba a NADA, en silencio (462 flechas en 61 guias: rutas del
% tipo «paleta → tipografia → lockup» salian rotas). Mapeamos el caracter a su
% version matematica, que si tiene glifo. Asi el autor puede seguir escribiendo
% «→» comodamente en el YAML.
\usepackage{newunicodechar}
\newunicodechar{→}{\ensuremath{\rightarrow}}
% Mismo problema con los simbolos de marca y vinetas: el «✓» de «una palomita ✓»
% salia como cuadrito vacio en el PDF de todas las guias que lo usaban.
\usepackage{pifont}
\newunicodechar{✓}{\ding{51}}
\newunicodechar{✗}{\ding{55}}
\newunicodechar{✔}{\ding{52}}
\newunicodechar{•}{\textbullet}
\newunicodechar{≥}{\ensuremath{\geq}}
\newunicodechar{≤}{\ensuremath{\leq}}

\setmainfont{Helvetica Neue}
\setsansfont{Helvetica Neue}
\setlength{\parindent}{0pt}
\setlength{\parskip}{6pt}
% Sin particion de palabras: en 6.º-7.º una palabra cortada por guion al final
% de linea («Comuni-cacion») frena la lectura mucho mas que una linea corta.
\hyphenpenalty=10000
\exhyphenpenalty=10000
\pretolerance=2000
\tolerance=3000
\setlength{\emergencystretch}{3em}
\linespread{1.10}
\renewcommand{\arraystretch}{1.25}
\setlist{leftmargin=5mm,itemsep=1mm,topsep=1mm}
\newcolumntype{Y}{>{\RaggedRight\arraybackslash}X}

\definecolor{milcMagenta}{HTML}{E6007E}
\definecolor{milcVino}{HTML}{5A0038}
\definecolor{milcTurquesa}{HTML}{0093A5}
\definecolor{milcVerde}{HTML}{58A923}
\definecolor{milcMostaza}{HTML}{E5B400}
\definecolor{milcOcre}{HTML}{B86600}
\definecolor{milcNegro}{HTML}{241816}
\definecolor{milcGris}{HTML}{F7F7F4}
\definecolor{milcLinea}{HTML}{E8E2DD}
\definecolor{milcGradePrimary}{HTML}{FF2D87}
\definecolor{milcGradeDark}{HTML}{9F1B56}
\definecolor{milcGradeSoft}{HTML}{FFE0EE}
\definecolor{milcGradeAccent}{HTML}{CFE7FF}
\definecolor{milcGradeText}{HTML}{FFFFFF}
\color{milcNegro}

\pagestyle{fancy}
\fancyhf{}
\lhead{\small Tecnología e Informática · Grado 10}
\rhead{\small Guía 26 · MILC v3}
\cfoot{\small\thepage}
\renewcommand{\headrulewidth}{0pt}

\newtcolorbox{titlebox}[1]{
  enhanced,colback=#1,colframe=#1,arc=7mm,boxrule=0pt,
  left=7mm,right=7mm,top=3mm,bottom=3mm,
  before skip=8pt,after skip=8pt,
  fontupper=\sffamily\bfseries\Large\color{white}
}

\newtcolorbox{softbox}[3]{
  enhanced,colback=#2,colframe=#1,borderline west={4pt}{0pt}{#1},
  arc=2mm,boxrule=.4pt,left=4mm,right=4mm,top=3mm,bottom=3mm,
  before skip=6pt,after skip=8pt,title={#3},coltitle=white,
  colbacktitle=#1,fonttitle=\sffamily\bfseries,fontupper=\RaggedRight,
  attach boxed title to top left={xshift=3mm,yshift=-2mm},
  boxed title style={arc=2mm, boxrule=0pt}
}

\newtcolorbox{drawbox}[2]{
  enhanced,colback=white,colframe=#1,arc=4mm,boxrule=1pt,
  left=5mm,right=5mm,top=4mm,bottom=4mm,before skip=8pt,after skip=8pt,
  title={#2},coltitle=white,colbacktitle=#1,fonttitle=\sffamily\bfseries,
  fontupper=\RaggedRight,
  attach boxed title to top left={xshift=4mm,yshift=-2mm},
  boxed title style={arc=3mm, boxrule=0pt}
}

\newtcolorbox{infoband}[1]{
  enhanced,colback=#1!8,colframe=#1!50,arc=2mm,boxrule=.5pt,
  left=4mm,right=4mm,top=2mm,bottom=2mm,before skip=4pt,after skip=4pt,
  fontupper=\small\RaggedRight
}

% Puente narrativo entre secciones (contrato editorial Conéctate).
% Da continuidad: "Acabas de X. Ahora vas a Y porque Z."
% Visualmente: banda izquierda en color del grado, texto en itálica pequeña.
\newtcolorbox{puentebox}{
  enhanced,colback=milcGradeSoft,colframe=milcGradeDark,
  borderline west={3pt}{0pt}{milcGradeDark},
  arc=0mm,boxrule=0pt,left=5mm,right=4mm,top=2.5mm,bottom=2.5mm,
  before skip=4pt,after skip=8pt,
  fontupper=\itshape\small\color{milcNegro}\RaggedRight
}
% La flecha va en modo matematico a proposito: Helvetica Neue NO tiene el glifo
% U+2192, asi que \textrightarrow se compilaba a NADA («Missing character: There
% is no → in font Helvetica Neue») y el puente salia sin flecha. En modo
% matematico la flecha viene de la fuente matematica y si se dibuja.
\newcommand{\puente}[1]{%
  \begin{puentebox}%
    \textcolor{milcGradeDark}{$\rightarrow$}\hspace{2mm}#1%
  \end{puentebox}%
}

\newcommand{\linead}[1]{\noindent\color{milcLinea}\rule{#1}{0.5pt}\color{milcNegro}}
\newcommand{\respuesta}[1]{\par\vspace{2mm}\foreach \n in {1,...,#1}{\noindent\color{milcLinea}\rule{\linewidth}{0.5pt}\par\vspace{5mm}}\color{milcNegro}}
\newcommand{\checkbox}{\textcolor{milcGradeDark}{\fbox{\phantom{x}}}\hspace{5pt}}

\newcommand{\phasecard}[4]{
\begin{tcolorbox}[enhanced,colback=#3,colframe=#2,arc=3mm,boxrule=.5pt,height=3.05cm,valign=center,halign=center,left=1.5mm,right=1.5mm,top=2mm,bottom=2mm]
{\sffamily\bfseries\fontsize{22}{24}\selectfont\textcolor{#2}{#1}}\par
{\sffamily\bfseries\small #4}
\end{tcolorbox}
}

% ── Piezas del rediseno 2026 (legibilidad 6.º-11.º) ───────────────────
% Banda de estacion: reemplaza el titlebox macizo. Titulo a la izquierda,
% dato practico (tiempo/modalidad) a la derecha, en una sola linea.
\newcommand{\milcestacion}[2]{%
\begin{tcolorbox}[enhanced,colback=#1,colframe=#1,arc=2mm,boxrule=0pt,
  left=5mm,right=5mm,top=2.6mm,bottom=2.6mm,before skip=11pt,after skip=7pt]
{\sffamily\bfseries\fontsize{15}{18}\selectfont\color{white}#2}
\end{tcolorbox}}

% Tarjeta de paso numerada: sustituye las tablas «Pilar / Aplicacion», que
% eran formato de consulta docente, no de estudiante. El numero va en un
% circulo sobre el borde para que la secuencia se lea de un vistazo.
\newcommand{\milcpaso}[3]{%
\begin{tcolorbox}[enhanced,colback=white,colframe=milcLinea,arc=1.5mm,boxrule=.9pt,
  left=14mm,right=4mm,top=3mm,bottom=3mm,before skip=4pt,after skip=4pt,
  fontupper=\RaggedRight,
  overlay={\node[anchor=north west,circle,fill=#2,text=white,inner sep=0pt,
    minimum size=7.4mm,font=\sffamily\bfseries\small]
    at ([xshift=3.6mm,yshift=-3.2mm]frame.north west) {#1};}]
#3
\end{tcolorbox}}

% Casilla marcable de tamano util con lapiz (la anterior era \fbox{\phantom{x}},
% de alto variable segun el texto que la seguia).
\newcommand{\milcmarca}{\framebox[3.8mm]{\rule{0pt}{3.4mm}}}

% Fila marcable de lista suelta (checks de la estacion 1).
\newcommand{\milccheckrow}[1]{%
\par\vspace{1.8mm}\noindent
\makebox[6.4mm][l]{\raisebox{-.55mm}{\textcolor{milcNegro}{\milcmarca}}}%
\parbox[t]{\dimexpr\linewidth-6.4mm\relax}{\RaggedRight #1}\par}

% Celda marcable dentro de la rubrica (tabla de autoevaluacion). Misma
% sangria francesa, pero pensada para vivir dentro de una columna X.
\newcommand{\milcopcion}[1]{%
\makebox[5.2mm][l]{\raisebox{-.55mm}{\textcolor{milcNegro}{\milcmarca}}}%
\parbox[t]{\dimexpr\linewidth-5.2mm\relax}{\RaggedRight #1}}

% ── Recursos visuales de la guía ──────────────────────────────────────
% Declarados en el bloque `recursos:` del YAML y emitidos por el builder.
% \guiaFigura   → imagen rasterizada o PDF (foto, carta celeste, montaje).
% \guiaDiagrama → fuente TikZ/circuitikz (.tex) incrustada como vector:
%                 es la vía para circuitos y esquemáticos editables.
% #1 = ruta relativa al .tex · #2 = pie (puede ir vacío)
\newcommand{\guiaFigura}[2]{%
\begin{center}
\includegraphics[width=0.88\linewidth,height=0.38\textheight,keepaspectratio]{#1}\\[1.5mm]
{\footnotesize\itshape\textcolor{milcNegro!75}{#2}}
\end{center}
}

\newcommand{\guiaDiagrama}[2]{%
\begin{center}
\input{#1}\\[1.5mm]
{\footnotesize\itshape\textcolor{milcNegro!75}{#2}}
\end{center}
}

\begin{document}
\thispagestyle{empty}

% ── Portada ───────────────────────────────────────────────────────────
% Antes: un rectangulo de color a pagina completa. Se veia bien en pantalla,
% pero en la fotocopiadora del colegio se comia el toner y salia gris sucio.
% Ahora la banda ocupa el tercio superior y el resto va en blanco.
\begin{tikzpicture}[remember picture,overlay]
  \path[fill=milcGradePrimary]
    (current page.north west) rectangle ([yshift=-10.6cm]current page.north east);

  \node[anchor=north west,text=milcGradeText] at ([xshift=2.0cm,yshift=-1.55cm]current page.north west)
    {\sffamily\bfseries\fontsize{12}{14}\selectfont GRADO};
  \node[anchor=north west,text=milcGradeText,opacity=.85] at ([xshift=2.0cm,yshift=-2.35cm]current page.north west)
    {\sffamily\bfseries\fontsize{8.5}{10}\selectfont SERIE GUÍAS · TECNOLOGÍA E INFORMÁTICA};
  \node[anchor=north east,text=milcGradeText,opacity=.85,align=right] at ([xshift=-2.0cm,yshift=-1.55cm]current page.north east)
    {\sffamily\bfseries\fontsize{9}{12}\selectfont I.E. Sor María Juliana\\Cartago, Valle del Cauca};

  \node[anchor=west,text=milcGradeText] (gradonum) at ([xshift=1.85cm,yshift=-7.1cm]current page.north west)
    {\sffamily\bfseries\fontsize{112}{112}\selectfont 10};
  % El circulito de grado (°) solo aplica a las guias de grado («11°»); en
  % Bebras y Semillero el numero grande es un momento/modulo, no un grado.
  % Se posiciona relativo al nodo `gradonum`, no en coordenadas absolutas.
  \draw[milcGradeText,line width=1.6pt]
    ([xshift=2.0mm,yshift=-4.5mm]gradonum.north east) circle (.15cm);

  \node[anchor=west,text=milcGradeText,text width=11.4cm,align=left] at ([xshift=1.5cm,yshift=0.5cm]gradonum.east)
    {
     {\sffamily\bfseries\fontsize{9}{11}\selectfont\MakeUppercase{Período 3 · Ofimática con IA: contabilidad, datos y emprendimiento}}\\[3mm]
     {\sffamily\bfseries\fontsize{21}{25}\selectfont\setlength{\baselineskip}{27pt}Encuestas y análisis\\[4pt]de mercado\\[4pt]con IA}\\[3.5mm]
     {\sffamily\bfseries\fontsize{15}{18}\selectfont Guía 26-10-TIC}
    };
\end{tikzpicture}

\vspace*{9.4cm}
\vfill

{\sffamily\bfseries\fontsize{8}{10}\selectfont\textcolor{milcGradeDark}{LO QUE TE LLEVAS HOY}}

\begin{tcolorbox}[enhanced,colback=white,colframe=milcNegro,arc=2mm,boxrule=1.6pt,
  left=5mm,right=5mm,top=4mm,bottom=4mm,before skip=3pt,after skip=12pt,
  fontupper=\RaggedRight\fontsize{11}{15.5}\selectfont]
Mini-estudio de mercado completo para un microemprendimiento real o hipotético (puede ser el mismo de la sesión 4) con (a) encuesta de 8-10 preguntas diseñada en Google Forms y revisada con IA (siguiendo método de la sesión 5 del P2), (b) aplicación a 15-20 personas del entorno cercano (familia, vecinos, compañeros, clientes potenciales), (c) análisis cualitativo y cuantitativo con asistencia de IA auditada, (d) 3 hallazgos numéricos clave, (e) 1 decisión razonada y documentada sobre qué producto o servicio ofrecer y con qué características
\end{tcolorbox}

\begin{tcolorbox}[enhanced,colback=milcGris,colframe=milcGris,arc=2mm,boxrule=0pt,
  left=5mm,right=5mm,top=3.5mm,bottom=3.5mm,after skip=12pt]
{\sffamily\bfseries\fontsize{8}{10}\selectfont\textcolor{milcNegro!55}{ESTA GUÍA ES DE}}\\[3mm]
\begin{tabularx}{\linewidth}{X p{3.2cm} p{3.2cm}}
\linead{\linewidth} & \linead{3.0cm} & \linead{3.0cm}\\[-1mm]
{\fontsize{8.5}{10}\selectfont\textcolor{milcNegro!55}{Nombre completo}} &
{\fontsize{8.5}{10}\selectfont\textcolor{milcNegro!55}{Grupo}} &
{\fontsize{8.5}{10}\selectfont\textcolor{milcNegro!55}{Fecha}}\\
\end{tabularx}
\end{tcolorbox}

\vfill

\noindent\textcolor{milcLinea}{\rule{\linewidth}{1pt}}\\[2mm]
{\sffamily\bfseries\fontsize{9.5}{12}\selectfont Autor: Dr. Álvaro Cárdenas Orozco}\\[1mm]
{\sffamily\fontsize{8.5}{11}\selectfont\textcolor{milcNegro!55}{Metodología MILC · apertura ancestral · pensamiento computacional · triángulo Dussel–estoicismo–Floridi}}

\newpage

\section*{Apertura: Encuestas y análisis de mercado con IA}

\begin{softbox}{milcGradeDark}{milcGradeSoft}{Saber ancestral}
En las tiendas, panaderías y talleres del centro de Cartago, hubo durante décadas una práctica que cualquier comerciante experimentado conocía: \textbf{nunca pedir mercancía sin antes preguntar a la clientela}. Antes de hacer pedido grande al proveedor, el tendero hablaba con sus clientes habituales: \emph{``¿qué les hace falta últimamente?''}, \emph{``¿qué les gustaría que trajera?''}, \emph{``¿qué precio pagarían?''}. La panadera preguntaba si querían más pan blanco o más integral. La modista preguntaba si la moda local iba hacia colores claros u oscuros. El tendero del barrio Obrero conocía a sus 50 clientes habituales por nombre y sabía qué preferían. Esa práctica tenía un nombre que los microempresarios no usaban pero practicaban: \textbf{investigación de mercado}. La sabiduría era ancestral y precisa: \textbf{el comerciante que no escucha a su clientela pierde la clientela}. La investigación de mercado moderna formaliza con IA lo que el tendero hacía de manera intuitiva: preguntar a la audiencia antes de invertir tiempo y dinero en un producto.
\end{softbox}

\begin{softbox}{milcTurquesa}{milcGris}{Saber contemporáneo}
Un \textbf{estudio de mercado} es la investigación que precede al lanzamiento de cualquier microemprendimiento. Responde 5 preguntas básicas: (1) \textbf{¿Quién es el cliente?}: edad, contexto, hábitos, ubicación. (2) \textbf{¿Qué necesita?}: problema real, no inventado. (3) \textbf{¿Cuánto está dispuesto a pagar?}: rango de precios aceptable. (4) \textbf{¿Cuál es la competencia?}: qué alternativas existen ya. (5) \textbf{¿Cómo nos enteramos?}: por qué canal el cliente conoce productos similares. La regla profesional: \textbf{ningún microemprendimiento sin estudio de mercado previo sobrevive a competencia básica}. El estudio combina: (a) \textbf{Encuesta cuantitativa}: Google Forms con preguntas cerradas a 15-20 personas. (b) \textbf{Análisis cualitativo}: respuestas abiertas analizadas con IA siguiendo método auditado de S6 del P2. (c) \textbf{Síntesis con hallazgos numéricos}: 3-5 cifras clave que dirijan decisión. (d) \textbf{Decisión razonada}: qué producto o servicio ofrecer, a qué precio, con qué características. La IA acelera análisis pero la decisión final sigue siendo del emprendedor.
\end{softbox}

\begin{softbox}{milcMostaza}{milcGris}{Pregunta-puente}
\textit{¿Qué sabía el tendero al preguntar a sus 50 clientes antes de hacer pedido grande, que el emprendedor novato olvida cuando lanza microempresa basado en \emph{``yo creo que la gente quiere esto''}? ¿Y por qué 15-20 entrevistados reales valen más que mil suposiciones?}
\end{softbox}

\begin{softbox}{milcVerde}{milcGris}{Saber y saber hacer}
Al terminar podrás: (1) \textbf{identificar} las 5 preguntas básicas que todo estudio de mercado debe responder; (2) \textbf{analizar} las respuestas de tu encuesta con asistencia de IA auditada; (3) \textbf{aplicar} las habilidades de S5 del P2 (encuesta con phronesis) al contexto comercial; (4) \textbf{evaluar} qué producto o servicio lanzar con criterio basado en datos reales.
\end{softbox}



\small
\begin{tabularx}{\textwidth}{>{\bfseries}p{3.0cm}Y}
\rowcolor{milcGradePrimary!12}Autor & Dr. Álvaro Cárdenas Orozco\\
DBA & Aplicación de ofimática asistida por IA a contabilidad básica y emprendimiento (MEN, T\&I)\\
Referentes & Lean Startup · Phronesis financiera · Ética del dato empresarial\\
Producto final & Mini-estudio de mercado completo para un microemprendimiento real o hipotético (puede ser el mismo de la sesión 4) con (a) encuesta de 8-10 preguntas diseñada en Google Forms y revisada con IA (siguiendo método de la sesión 5 del P2), (b) aplicación a 15-20 personas del entorno cercano (familia, vecinos, compañeros, clientes potenciales), (c) análisis cualitativo y cuantitativo con asistencia de IA auditada, (d) 3 hallazgos numéricos clave, (e) 1 decisión razonada y documentada sobre qué producto o servicio ofrecer y con qué características\\
\end{tabularx}
\normalsize

\puente{Antes de empezar, mira el viaje completo. En la sesión 4 proyectaste flujo de caja de un microemprendimiento. Hoy validas si tu microemprendimiento tiene mercado real. Las próximas sesiones (Canvas, comunicación, proyecto integrador, sustentación) usarán los hallazgos del estudio de mercado.}

\milcestacion{milcGradeDark}{Ruta MILC: así vamos a aprender}
\begin{tabularx}{\textwidth}{>{\centering\arraybackslash}X>{\centering\arraybackslash}X>{\centering\arraybackslash}X>{\centering\arraybackslash}X}
\phasecard{1}{milcVerde}{milcGris}{Escucha\\[-1mm]\small Observo, escucho y pregunto.} &
\phasecard{2}{milcTurquesa}{milcGris}{Entender\\[-1mm]\small Sistematizo y ordeno.} &
\phasecard{3}{milcMagenta}{milcGris}{Hacer\\[-1mm]\small Construyo y aplico.} &
\phasecard{4}{milcVino}{milcGris}{Cierre\\[-1mm]\small Evalúo y reflexiono.}
\end{tabularx}


\puente{Antes de teorizar, vas a hacer una \emph{conversación corta}: piensa en el microemprendimiento de la sesión 4 (real o hipotético) y conversa 5 minutos con 1 persona de tu entorno (mamá, hermano, vecino) sobre si compraría el producto o servicio. Anota qué te dijo.}

\milcestacion{milcVerde}{Estación 1 de 3 · Escucha}

\begin{softbox}{milcVerde}{milcGris}{Escena para escuchar}
\textbf{Actividad 1 · IDENTIFICA --- Conversación corta con 1 cliente potencial} (15 min · individual). Toma el microemprendimiento de la sesión 4. Conversa 5-10 minutos con 1 persona de tu entorno (mamá, hermano, vecino, amigo) que podría ser cliente potencial. Pregunta de manera natural: (1) ¿usarías un producto como X? (2) ¿qué pagarías por él? (3) ¿qué te detendría de comprarlo? (4) ¿qué alternativas conoces ya? (5) ¿cómo te enteras de productos similares? Anota textualmente lo que te dijo (en sus palabras, no tu interpretación).
\end{softbox}

{\sffamily\bfseries\fontsize{8}{10}\selectfont\textcolor{milcNegro!55}{MARCA CUANDO LO HAYAS HECHO}}
\milccheckrow{Tuve conversación de 5-10 min con 1 persona.}
\milccheckrow{Hice las 5 preguntas con naturalidad.}
\milccheckrow{Anoté las respuestas textuales.}

\begin{infoband}{milcVerde}


\textbf{Lo que va al cuaderno (Actividad 1 · IDENTIFICA).} Título: «Actividad 1 --- Conversación 1 cliente». \textbf{Formato:} ficha con las 5 preguntas y respuestas textuales. \textbf{Sabes que terminaste cuando} tienes voz real de un cliente potencial.
\end{infoband}

\puente{Ya tienes 1 conversación. Ahora vas a entender la \textbf{anatomía del estudio de mercado}: 5 preguntas básicas, estructura de la encuesta, análisis con IA.}

\milcestacion{milcTurquesa}{Estación 2 de 3 · Entender}

\begin{softbox}{milcTurquesa}{milcGris}{Lo que vas a entender}
Tu conversación te dio voz real. Ahora necesitas la \textbf{anatomía profesional} del estudio de mercado con asistencia de IA.
\end{softbox}

\milcpaso{1}{milcTurquesa}{Las \textbf{5 preguntas básicas} del estudio de mercado se descomponen así: (1) \textbf{Cliente}: ¿quién es la persona típica que compraría? Edad, contexto, hábitos. (2) \textbf{Necesidad}: ¿qué problema real resuelve mi producto? No invento; verifico. (3) \textbf{Precio}: ¿cuánto está dispuesto a pagar? Rango de precios aceptable. (4) \textbf{Competencia}: ¿qué alternativas existen? (otros productos, marcas, formas de resolver el mismo problema). (5) \textbf{Canal}: ¿cómo se entera el cliente de productos similares? (redes, voz a voz, publicidad).}
\milcpaso{2}{milcTurquesa}{La \textbf{estructura de la encuesta de mercado} en Google Forms: (a) \textbf{Sección demográfica} (2-3 preguntas): edad, contexto, hábito relacionado. (b) \textbf{Sección de necesidad} (2-3 preguntas): ¿enfrentas el problema X?, ¿con qué frecuencia? (c) \textbf{Sección de oferta} (2-3 preguntas): ¿usarías un producto como Y?, ¿qué pagarías? (d) \textbf{Sección abierta} (1-2 preguntas): ¿qué te detendría de comprar?, ¿qué mejorarías? Total: 8-10 preguntas. Tiempo de respuesta: 5-7 minutos.}
\milcpaso{3}{milcTurquesa}{El \textbf{análisis con IA auditada} sigue el método de S6 del P2: (1) recolectar 15-20 respuestas. (2) pasar respuestas abiertas a ChatGPT/Gemini con prompt de clasificación. (3) auditar respuesta por respuesta. (4) sacar hallazgos numéricos: \emph{``17 de 20 personas dijeron que pagarían entre 5.000 y 8.000''}, \emph{``12 mencionaron que les gustaría versión vegetariana''}. La regla: \textbf{los hallazgos son cifras, no opiniones del emprendedor}.}
\milcpaso{4}{milcTurquesa}{Algoritmo: (1) tomar microemprendimiento de S4. (2) diseñar encuesta 8-10 preguntas siguiendo S5 del P2 (con filtro de IA contra sesgos). (3) aplicar a 15-20 personas durante la semana. (4) analizar con IA auditada. (5) sacar 3 hallazgos numéricos clave. (6) decidir qué ajustar en tu microemprendimiento.}

\begin{drawbox}{milcTurquesa}{Anatomía del estudio de mercado}
\textbf{5 preguntas básicas:} cliente / necesidad / precio / competencia / canal. \\[1mm] \textbf{Encuesta:} 8-10 preguntas en 4 secciones. \\[1mm] \textbf{Aplicación:} 15-20 personas del entorno. \\[1mm] \textbf{Análisis:} IA auditada como en S6 del P2. \\[1mm] \textbf{Hallazgos:} 3-5 cifras clave para decisión.
\end{drawbox}

\begin{softbox}{milcMostaza}{milcGris}{Errores comunes y su corrección}
(1) \textbf{Encuesta a solo amigos cercanos}: dirán que sí por cortesía. Incluir 3+ desconocidos. (2) \textbf{Preguntas hipotéticas sin precio}: \emph{``¿comprarías?''} sin decir cuánto cobra produce sí falsos. (3) \textbf{No analizar respuestas abiertas}: solo cuantitativo pierde matices. (4) \textbf{Aceptar lo que la IA dice sin auditar}. (5) \textbf{Decisión sin sustento numérico}: \emph{``creo que vendería bien''} no es decisión.
\end{softbox}

\begin{infoband}{milcTurquesa}


\textbf{Lo que va al cuaderno (Actividad 2 · ANALIZA).} Título: «Actividad 2 --- Anatomía estudio mercado». \textbf{Formato:} (a) las 5 preguntas básicas; (b) estructura encuesta; (c) los 5 errores con marca. \textbf{Sabes que terminaste cuando} sabes diseñar encuesta de mercado sin abrir esta guía.
\end{infoband}

\puente{Tienes el método. Toca aplicarlo: diseñas encuesta de 8-10 preguntas, la aplicas a 15-20 personas durante la semana, analizas con IA auditada, sacas 3 hallazgos, decides.}

\milcestacion{milcMagenta}{Estación 3 de 3 · Hacer}

\begin{softbox}{milcMagenta}{milcGris}{Cómo vas a construir tu producto}
\textbf{Actividad 3 · APLICA + EVALÚA --- Estudio de mercado completo} (40 min en sesión + semana de aplicación · individual). Hora de aplicar. Diseñas encuesta, la aplicas a 15-20 personas, analizas con IA auditada, sacas hallazgos, decides.
\end{softbox}

\milcpaso{1}{milcMagenta}{Tu estudio se construye en 7 pasos: (1) tomar microemprendimiento de S4. (2) diseñar 8-10 preguntas en Forms cubriendo las 5 preguntas básicas. (3) revisar con filtro de IA para sesgos (S5 del P2). (4) durante la semana, aplicar a 15-20 personas. Mezclar: 5 cercanos, 5 menos cercanos, 5 desconocidos cuando se pueda. (5) analizar respuestas con IA + auditoría. (6) sacar 3 hallazgos numéricos clave. (7) decidir cambios al microemprendimiento basados en hallazgos.}
\milcpaso{2}{milcMagenta}{Sigue el patrón \textbf{``pivotar con datos''}: el objetivo del estudio no es confirmar tu idea, es \textbf{descubrir si necesita ajuste}. Si los hallazgos indican que el precio que pensabas era muy alto, baja el precio. Si la mayoría no enfrenta el problema, replantea audiencia o producto. La phronesis del emprendedor consiste en \textbf{escuchar los datos aunque contradigan la idea inicial}.}
\milcpaso{3}{milcMagenta}{Los \textbf{3 hallazgos numéricos} se formulan como afirmaciones con cifra. Ejemplos: (a) \emph{``17 de 20 encuestados (85 por ciento) enfrentan el problema X al menos 1 vez por semana''}. (b) \emph{``El precio promedio que están dispuestos a pagar es 6.500 pesos''}. (c) \emph{``12 de 20 (60 por ciento) prefieren versión sin gluten''}. Estos hallazgos dirigen decisiones.}
\milcpaso{4}{milcMagenta}{Algoritmo: (1) diseñar encuesta. (2) filtro IA contra sesgos. (3) aplicar a 15-20 personas. (4) análisis con IA + auditoría. (5) 3 hallazgos numéricos. (6) decisión documentada de ajustes al microemprendimiento.}

\begin{drawbox}{milcMagenta}{Checklist al cerrar el estudio}
\checkbox Encuesta de 8-10 preguntas diseñada. \\[1mm] \checkbox Filtro IA contra sesgos aplicado. \\[1mm] \checkbox 15-20 respuestas recolectadas. \\[1mm] \checkbox Mezcla de cercanos, menos cercanos y desconocidos. \\[1mm] \checkbox Análisis con IA auditado. \\[1mm] \checkbox 3 hallazgos numéricos claros. \\[1mm] \checkbox Decisión documentada de ajustes.
\end{drawbox}

\begin{softbox}{milcMagenta}{milcGris}{Plantilla de trabajo}
\textbf{Microemprendimiento.} \textit{Real o hipotético, de S4.} \\[1mm] \textbf{Encuesta.} \textit{8-10 preguntas en 4 secciones.} \\[1mm] \textbf{Aplicación.} \textit{15-20 personas mixtas.} \\[1mm] \textbf{Análisis.} \textit{IA + auditoría humana.} \\[1mm] \textbf{Hallazgos.} \textit{3 afirmaciones con cifra.} \\[1mm] \textbf{Decisión.} \textit{Ajustes documentados.}
\end{softbox}

\begin{infoband}{milcMagenta}


\textbf{Lo que va al cuaderno (Actividad 3 · APLICA + EVALÚA).} Título: «Actividad 3 --- Estudio de mercado». \textbf{Formato:} (a) link encuesta Forms; (b) tabla resultados; (c) análisis auditado; (d) 3 hallazgos; (e) decisión justificada. \textbf{Sabes que terminaste cuando} sabes con datos si tu microemprendimiento tiene mercado.
\end{infoband}

\puente{Lo que sigue resume \emph{qué} entregas hoy: estudio de mercado completo + hallazgos + decisión. El criterio principal: que tu decisión final tenga sustento en datos, no en intuición.}

\milcestacion{milcMostaza}{Producto de la sesión}
\textbf{Estudio mercado + 15-20 respuestas + análisis IA + 3 hallazgos + decisión}.
\begin{softbox}{milcMostaza}{milcGris}{Criterios de calidad}
(1) Encuesta 8-10 preguntas con filtro IA. (2) 15-20 respuestas con mezcla de audiencias. (3) Análisis IA auditado. (4) 3 hallazgos numéricos. (5) Decisión documentada con ajustes. (6) Microemprendimiento alineado con datos.
\end{softbox}

\puente{Antes de cerrar, mira el estudio desde las cinco dimensiones humanas. Preguntar a la audiencia antes de lanzar es respeto por el cliente potencial; inventar lo que quiere es desperdicio garantizado.}

\milcestacion{milcVino}{Evaluación liberadora · cinco dimensiones}

\begin{softbox}{milcVino}{milcGris}{Sentido de la evaluación}
Evaluar no es solo revisar una nota. Es reconocer qué aprendí, qué cambió en mí, qué emociones manejé, cómo me relaciono con la ciudad y la comunidad, y qué le dejaré a quien viene detrás.
\end{softbox}

\textbf{Mi reflexión liberadora en cinco dimensiones.} Completa cada frase iniciada:

\small
\begin{tabularx}{\textwidth}{>{\bfseries\RaggedRight\arraybackslash}p{3.4cm}Y>{\RaggedRight\arraybackslash}p{4.6cm}}
\rowcolor{milcVino}\color{white}Dimensión & \color{white}\bfseries Mi respuesta (frame starter) & \color{white}\bfseries Algo que descubrí\\
1. Personal & Lo que más cambió en mí fue \linead{6.5cm} & \linead{4.2cm}\\
2. Emocional & La emoción más fuerte fue \linead{2.5cm} y la regulé \linead{3cm} & \linead{4.2cm}\\
3. Ciudadana & En democracia esto importa porque \linead{6cm} & \linead{4.2cm}\\
4. Local (Cartago) & En mi barrio sería útil para \linead{6.5cm} & \linead{4.2cm}\\
5. Intergeneracional & A mi abuela/abuelo le diría \linead{6.5cm} & \linead{4.2cm}\\
\end{tabularx}
\normalsize

\begin{infoband}{milcVino}
\textbf{Para la próxima sustentación.} Vuelve a esta tabla en seis meses. Verás que las respuestas cambian — no porque lo aprendido cambie, sino porque tú cambias. La evaluación liberadora es una conversación contigo a través del tiempo.
\end{infoband}


\puente{Y para terminar, tres voces sobre la escucha al mercado. Dussel sobre los estudios que escuchan a comunidades reales, Marco Aurelio sobre la disciplina de preguntar antes de afirmar, Floridi sobre el dato del cliente como ética emprendedora.}

\milcestacion{milcGradeDark}{Triángulo de pensamiento · cierre filosófico}

\begin{softbox}{milcVino}{milcGris}{Dussel · lente del nosotros}
\textit{Escuchar al cliente potencial antes de lanzar es respeto por la comunidad que se pretende servir.}

{\footnotesize\itshape\color{milcNegro!70}--- formulación de esta guía, al modo de Enrique Dussel. No es una cita textual.}

Cuando preguntas a 15-20 personas reales antes de lanzar microempresa, estás respetando a esa comunidad como sujeto, no objeto de venta. La filosofía de la liberación pide ese acto: \textbf{escuchar antes de ofrecer}. La phronesis del emprendedor consiste en \textbf{ajustar el producto a la voz del cliente, no imponer lo que el emprendedor quiere vender}.

\textbf{¿Estoy escuchando a mis clientes potenciales antes de lanzar, o impongo lo que yo creo que quieren?}
\end{softbox}
\linead{\linewidth}\\[1mm]
\linead{\linewidth}

\begin{softbox}{milcMostaza}{milcGris}{Estoicismo · Marco Aurelio · cuidado interior}
\textit{Preguntar antes de afirmar es virtud emprendedora; lanzar sin estudio es vanidad que se paga caro.}

{\footnotesize\itshape\color{milcNegro!70}--- formulación de esta guía, al modo de Marco Aurelio. No es una cita textual.}

Es tentador saltarse el estudio porque \emph{``ya sé lo que la gente quiere''}. La sabiduría estoica recomienda lo opuesto: \emph{pregunta antes de afirmar}. La intuición del emprendedor a veces acierta y muchas veces no. El estudio convierte intuición en información verificable.

\textbf{¿Estoy haciendo estudio honesto, o me lo salto creyendo que ya sé la respuesta?}
\end{softbox}
\linead{\linewidth}\\[1mm]
\linead{\linewidth}

\begin{softbox}{milcTurquesa}{milcGris}{Floridi · lente de la infoesfera}
\textit{Los datos del cliente bien recolectados son ética emprendedora en la era de la información distribuida.}

{\footnotesize\itshape\color{milcNegro!70}--- formulación de esta guía, al modo de Luciano Floridi. No es una cita textual.}

En la era de la información, los emprendedores que tienen datos sobre sus clientes ganan ventaja sobre los que no. La ética emprendedora exige recolectar esos datos con respeto: consentimiento informado, preguntas neutras, análisis honesto. Tu estudio de hoy te entrena en esa práctica que rige todo emprendimiento serio del siglo XXI.

\textbf{¿Mis datos del cliente están recolectados con rigor profesional o solo con encuestas casuales?}
\end{softbox}
\linead{\linewidth}\\[1mm]
\linead{\linewidth}

\begin{infoband}{milcNegro}
\textbf{Nota para el docente.} Las tres formulaciones son del autor de la guía, no citas textuales de los pensadores: recogen su lente para pensar el tema, no sus palabras. Si vas a citarlos en clase, remite a la obra original.
\end{infoband}

\milcestacion{milcVino}{Mi autoevaluación · marca una casilla por fila}

{\small
\setlength{\extrarowheight}{2.2pt}
\begin{tabularx}{\textwidth}{>{\RaggedRight\arraybackslash\bfseries}p{3.0cm}YYY}
\rowcolor{milcVino}\color{white}\bfseries Criterio & \color{white}\bfseries Lo logré & \color{white}\bfseries En proceso & \color{white}\bfseries Necesito apoyo\\[.6mm]
Comprensión del tema & \milcopcion{Explico las ideas con mis palabras.} & \milcopcion{Reconozco las ideas, falta precisión.} & \milcopcion{Necesito repasar lo básico.}\\[.8mm]
\rowcolor{milcGris}Pensamiento computacional & \milcopcion{Usé los pilares al armar mi producto.} & \milcopcion{Usé algunos pilares.} & \milcopcion{Necesito releer la estación 2.}\\[.8mm]
Aplicación y producto & \milcopcion{Mi producto cumple los criterios.} & \milcopcion{Está armado, con ajustes pendientes.} & \milcopcion{Aún está en construcción.}\\[.8mm]
\rowcolor{milcGris}Reflexión 5 dimensiones & \milcopcion{Respondí cada una con honestidad.} & \milcopcion{Respondí algunas con detalle.} & \milcopcion{Mis respuestas son muy generales.}\\[.8mm]
Triángulo de pensamiento & \milcopcion{Conecté las 3 voces con mi trabajo.} & \milcopcion{Conecté al menos una.} & \milcopcion{No las relaciono con mi proceso.}\\[.8mm]
\end{tabularx}}

\vspace{3mm}
\textbf{Hoy entendí que:}\\[1mm]
\linead{\linewidth}\\[1mm]
\linead{\linewidth}

\textbf{Mi compromiso para la próxima sesión es:}\\[1mm]
\linead{\linewidth}\\[1mm]
\linead{\linewidth}

\begin{softbox}{milcMostaza}{milcGris}{Cita de despedida · Marco Aurelio}
\textit{``El obstáculo en el camino se vuelve el camino. Lo que estorba al camino se vuelve camino.''}\\[1mm]
Cada error de hoy, bien observado, se convierte en pieza del camino que pisarás mañana.
\end{softbox}

\small
\begin{tabularx}{\textwidth}{>{\bfseries}p{3.6cm}Y}
\rowcolor{milcGradeDark!12}Nombre completo: & \linead{10cm}\\
Fecha: & \linead{4cm} \quad Curso/grupo: \linead{4cm}\\
Mi firma: & \linead{9cm}\\
\end{tabularx}
\normalsize

\begin{infoband}{milcGradeDark}
\textbf{Para la próxima sesión.} Llega con tu estudio de mercado completo y los 3 hallazgos. En la sesión 7 vas a usar esos datos para construir el \textbf{Business Model Canvas} con IA: el lienzo del negocio en 9 casillas. El estudio de hoy es la voz del mercado; el Canvas de mañana es la arquitectura del modelo.
\end{infoband}

\end{document}
